% ============================================================
\section{Phần II: Diễn giải ngữ nghĩa và hỏi đáp}
\label{sec:part2}
% ============================================================

\subsection{Biểu diễn ngữ nghĩa}

Mỗi câu được parse thành công được diễn giải thành một \textbf{frame ngữ nghĩa}
\code{SemanticFrame}, ghi ra \code{semantic.txt} dưới dạng vị từ
\texttt{INTENT(slot=giá trị, ...)}:
\begin{itemize}
  \item \code{intent}: loại yêu cầu (\cref{tab:branches});
  \item \code{entity}: thực thể chính đã chuẩn hóa về ID của KB, ví dụ \code{WEEK_03},
        \code{CH03}, \code{LO2.5}, \code{PART_II}, \code{CO3085};
  \item \code{slots}: các tham số bổ sung như \code{field=CREDITS},
        \code{relation=PARSING}, \code{feature=parse tree};
  \item \code{source_branch}: nhánh \texttt{Q\_*} đã sinh ra frame, giữ lại để giải
        thích kết quả.
\end{itemize}
File \code{intent-entity.txt} ghi cùng thông tin dưới dạng JSON Lines, mỗi câu một
dòng gồm câu hỏi, intent gốc, intent sau khi giải ngữ cảnh và entity:
\begin{trace}
{"entity": "WEEK_03", "id": 4, "intent": "GET_SCHEDULE", "question": "Tuần 3 học gì?", "resolved_intent": "GET_SCHEDULE"}
\end{trace}

\subsection{Từ cây cú pháp sang frame}

Bộ diễn giải (\code{SemanticInterpreter}) là một tập luật ánh xạ cây sang frame theo
nguyên tắc \emph{tính hợp thành}: nghĩa của câu được ghép từ nghĩa của các thành phần có
tên trên cây.
\begin{enumerate}
  \item Tìm node con duy nhất của \texttt{QUESTION}; tên nhánh \texttt{Q\_*} xác định
        intent và chọn \emph{handler} tương ứng.
  \item Handler duyệt cây con để lấy các thành phần có tên (\texttt{WEEK\_ENTITY},
        \texttt{CHAPTER\_ENTITY}, \texttt{TOPIC\_ENTITY}, \texttt{RULE\_TOPIC},
        \texttt{ASSIGNMENT\_FEATURE}, \dots) và cụm từ mà chúng bao phủ.
  \item Entity lexicon (\code{entities.py}, xây từ \code{entities.txt} và bảng alias)
        chuẩn hóa cụm từ về ID: ``tuần 3'' $\rightarrow$ \code{WEEK_03},
        ``NLP'' $\rightarrow$ \code{CO3085}, ``thư viện có sẵn'' $\rightarrow$
        \code{EXTERNAL_LIBRARIES}.
  \item Handler điền các slot còn lại từ nhãn trên cây, ví dụ node
        \texttt{FEAT\_PARSING} cho biết câu hỏi về Phần I của bài tập lớn.
\end{enumerate}
Bộ diễn giải \textbf{không truy cập KB}: nó chỉ trả lời câu hỏi ``người dùng muốn
gì'', còn việc dữ kiện có tồn tại hay không thuộc về tầng truy vấn.

\begin{table}[htbp]
  \centering
  \small
  \begin{tabularx}{\linewidth}{|P{4cm}|L|}
    \hline
    \textbf{Câu hỏi} & \textbf{Frame ngữ nghĩa} \\ \hline
    Môn NLP có bao nhiêu tín chỉ? & \code{GET_COURSE_INFO(entity=CO3085, course=CO3085, field=CREDITS)} \\ \hline
    Tuần 3 học gì? & \code{GET_SCHEDULE(entity=WEEK_03, week=WEEK_03)} \\ \hline
    BTL có phần parse tree không? & \code{GET_ASSIGNMENT(entity=PART_I, assignment=BTL01, section=PART_I, feature=parse tree)} \\ \hline
    Có được dùng thư viện có sẵn không? & \code{GET_RULE(entity=EXTERNAL_LIBRARIES, rule=EXTERNAL_LIBRARIES)} \\ \hline
    Ai dạy NLP? & \code{GET_INSTRUCTOR(entity=CO3085, course=CO3085)} \\ \hline
    Nó có liên quan đến parsing không? & \code{CONTEXT_DEPENDENT(action=TOPIC_RELATION, relation=PARSING)} \\ \hline
    Hôm nay trời có mưa không? & \code{UNKNOWN()} -- không có cây \\ \hline
  \end{tabularx}
  \caption{Ví dụ frame ngữ nghĩa (trích từ kết quả chạy thật)}
  \label{tab:frames}
\end{table}

\subsection{Truy vấn knowledge base}

Bộ truy vấn (\code{QueryEngine}) chọn cách tra cứu theo intent và slot, rồi trả về
\code{QueryResult} gồm trạng thái \code{found}, loại dữ kiện \code{kind}, dữ liệu, danh
sách file nguồn và lý do khi không tìm thấy. Biểu thức truy vấn được ghi ra
\code{query.txt} để người đọc thấy hệ thống đã tra ở đâu (\cref{tab:queries}).

\begin{table}[htbp]
  \centering
  \small
  \begin{tabularx}{\linewidth}{|P{3.3cm}|P{6.2cm}|L|}
    \hline
    \textbf{Intent} & \textbf{Biểu thức truy vấn} & \textbf{Cách tra} \\ \hline
    \code{GET_COURSE_INFO} & \code{KB.course[CREDITS]} & Theo trường của môn học \\ \hline
    \code{GET_SCHEDULE} & \code{KB.weeks[WEEK_03]}, \code{KB.chapter_weeks[CH04]} & Theo ID tuần hoặc chỉ mục chương--tuần \\ \hline
    \code{GET_TOPIC} & \code{KB.topics[CH03]}, \code{KB.topics[CH03].contains(PARSING)} & Theo ID chương; tìm chủ đề theo từ khóa \\ \hline
    \code{GET_LO} & \code{KB.learning_outcomes[LO2.5]} & Theo ID chuẩn đầu ra \\ \hline
    \code{GET_ASSIGNMENT} & \code{KB.assignment.sections[PART_I]} & Theo phần của bài tập lớn \\ \hline
    \code{GET_DEADLINE} & \code{KB.weeks[*].milestone} & Gom các mốc trong lịch học \\ \hline
    \code{GET_RESOURCE} & \code{KB.resources[...]} & Chấm điểm tên và mô tả tài liệu \\ \hline
    \code{GET_RULE} & \code{KB.regulations[external_libraries]} & Theo khóa quy định \\ \hline
    \code{GET_INSTRUCTOR} & \code{KB.course[instructor]} & KB không có, luôn \code{NOT_FOUND} \\ \hline
  \end{tabularx}
  \caption{Cách truy vấn theo intent}
  \label{tab:queries}
\end{table}

Tra cứu theo ID (tuần, chương, LO, quy định) là tra cứu chính xác trên collection đã
nạp. Tìm chủ đề ưu tiên từ khóa khớp nguyên cụm rồi mới xét nội dung; tìm tài liệu chấm
điểm theo tên và mô tả cách sử dụng. Quan hệ giữa LO và chương dùng bảng ánh xạ cố định
của miền này, vì KB không lưu quan hệ đó.

\subsection{Sinh câu trả lời và xử lý trường hợp không trả lời được}

Bộ sinh câu trả lời chọn template tiếng Việt theo \code{kind} và điền dữ kiện từ
\code{QueryResult}, ví dụ template \code{SCHEDULE_WEEK} nêu tuần, chương, chủ đề và mốc
bài tập lớn nếu có. Bất biến quan trọng: một truy vấn \code{found=True} luôn cho câu trả
lời chứa dữ kiện; ngược lại câu trả lời rơi vào một trong ba fallback, tùy nguyên nhân:

\begin{table}[htbp]
  \centering
  \small
  \begin{tabularx}{\linewidth}{|P{3.6cm}|P{3.4cm}|L|}
    \hline
    \textbf{Nguyên nhân} & \textbf{Dấu hiệu} & \textbf{Câu trả lời} \\ \hline
    Câu ngoài văn phạm hoặc ngoài miền & Cây \texttt{()}, intent \code{UNKNOWN} &
      Xin lỗi, câu hỏi không thuộc văn phạm hoặc ngoài phạm vi Course Assistant. \\ \hline
    Thiếu ngữ cảnh & \code{CONTEXT_DEPENDENT} chưa giải được &
      Xin lỗi, tôi chưa có đủ ngữ cảnh để hiểu tham chiếu trong câu hỏi này. \\ \hline
    KB thiếu dữ kiện & Có cây và intent, \code{found=False} &
      Xin lỗi, tôi không tìm thấy thông tin này trong dữ liệu của môn học. \\ \hline
  \end{tabularx}
  \caption{Ba loại fallback}
  \label{tab:fallback}
\end{table}

Câu thứ ba là nguyên văn thông báo mà đề yêu cầu \cite{spec}. Việc tách ba trường hợp
giúp người dùng biết nên diễn đạt lại câu hỏi, bổ sung ngữ cảnh, hay dữ kiện đó thực sự
không có trong KB.

\subsection{Ví dụ trace đầy đủ}

Chế độ \code{--trace} in toàn bộ các bước cho một câu. Hai ví dụ dưới đây lấy từ lần
chạy thật: câu trả lời được, và câu hiểu được nhưng KB không có dữ kiện.
\begin{trace}
INPUT: Môn NLP có bao nhiêu tín chỉ?
TOKENS: môn / nlp / có / bao / nhiêu / tín / chỉ / ?
PARSE: (S (PREFIX ε) (QUESTION (Q_COURSE_CREDITS (COURSE_SUBJECT môn (COURSE_ENTITY nlp)) (CREDIT_VERB có) (CREDIT_QUESTION bao nhiêu tín chỉ))) (SUFFIX (END_MARK ?)))
SEMANTIC: GET_COURSE_INFO(entity=CO3085, course=CO3085, field=CREDITS)
QUERY: KB.course[CREDITS]
STATUS: FOUND
SOURCE: kb/course_info.txt
ANSWER: Môn Xử lý Ngôn ngữ Tự nhiên (CO3085) có 3 tín chỉ.
\end{trace}
\begin{trace}
INPUT: Ai dạy NLP?
PARSE: (S (PREFIX ε) (QUESTION (Q_INSTRUCTOR ai dạy (COURSE_SUBJECT (COURSE_ENTITY nlp)))) (SUFFIX (END_MARK ?)))
SEMANTIC: GET_INSTRUCTOR(entity=CO3085, course=CO3085)
QUERY: KB.course[instructor]
STATUS: NOT_FOUND
ANSWER: Xin lỗi, tôi không tìm thấy thông tin này trong dữ liệu của môn học.
\end{trace}
